\chapter{Chapter06 — Fourier To F₂Polynomial}
%%% generated by the blueprint pipeline from FABL.Chapter06.F₂Polynomials.FourierToF₂Polynomial
\section{Overview}

Declarations of the Lean module \texttt{FABL.Chapter06.F₂Polynomials.FourierToF₂Polynomial}.

\section{Declarations}

\begin{definition}[Support of numerical coefficients]
\lean{FABL.numericalSupport}
\statementsource{Analysis-of-Boolean-Functions-by-Ryan-ODonnell}{
  pdf-ffbe9c7b3306-p018-b002,
  pdf-ffbe9c7b3306-p025-b004,
  pdf-ffbe9c7b3306-p150-b001
}
\label{FABL.numericalSupport}
\leanok
\uses{FABL.NumericalCoefficients}
For a family of real coefficients $c(S)$ indexed by subsets $S$ of $[n]$, define its
support to be the collection of subsets $S\\subseteq[n]$ for which $c(S)\\ne0$.
\end{definition}

\begin{theorem}[Membership in numerical support]
\lean{FABL.mem_numericalSupport}
\label{FABL.mem_numericalSupport}
\leanok
\uses{FABL.NumericalCoefficients, FABL.numericalSupport}
\difficulty{1}
Let $c$ assign a real number to each subset of $\{0,\ldots,n-1\}$. A subset $S$ belongs
to the numerical support of $c$ if and only if $c(S)\ne 0$.
\end{theorem}

\begin{definition}[Degree of numerical coefficients]
\lean{FABL.numericalDegree}
\statementsource{Analysis-of-Boolean-Functions-by-Ryan-ODonnell}{
  pdf-ffbe9c7b3306-p036-b002,
  pdf-ffbe9c7b3306-p149-b004,
  pdf-ffbe9c7b3306-p150-b001
}
\label{FABL.numericalDegree}
\leanok
\uses{FABL.NumericalCoefficients, FABL.numericalSupport}
For a real coefficient family $c(S)$ indexed by subsets of $[n]$, define its numerical
degree to be the largest cardinality of a subset $S$ with nonzero coefficient, taking
the value $0$ when all coefficients vanish.
\end{definition}

\begin{theorem}[Characterization of bounded numerical degree]
\lean{FABL.numericalDegree_le_iff}
\proofsource{Analysis-of-Boolean-Functions-by-Ryan-ODonnell}{
  pdf-ffbe9c7b3306-p014-b007,
  pdf-ffbe9c7b3306-p036-b002
}
\label{FABL.numericalDegree_le_iff}
\leanok
\uses{FABL.NumericalCoefficients, FABL.numericalDegree}
\difficulty{4}
\proofstep
  {FABL.𝔽₂}
  {direct}
  {Analysis-of-Boolean-Functions-by-Ryan-ODonnell}
  {pdf-ffbe9c7b3306-p015-b008}
\proofstep
  {FABL.F₂Cube}
  {context}
  {Analysis-of-Boolean-Functions-by-Ryan-ODonnell}
  {pdf-ffbe9c7b3306-p023-b001,
    pdf-ffbe9c7b3306-p148-b003}
\proofstep
  {FABL.Sign}
  {context}
  {Analysis-of-Boolean-Functions-by-Ryan-ODonnell}
  {pdf-ffbe9c7b3306-p014-b002,
    pdf-ffbe9c7b3306-p022-b004}
\proofstep
  {FABL.signEncode}
  {direct}
  {Analysis-of-Boolean-Functions-by-Ryan-ODonnell}
  {pdf-ffbe9c7b3306-p022-b004}
\proofstep
  {FABL.binarySignEquiv}
  {context}
  {Analysis-of-Boolean-Functions-by-Ryan-ODonnell}
  {pdf-ffbe9c7b3306-p022-b004}
\proofstep
  {FABL.binaryCubeSignEquiv}
  {context}
  {Analysis-of-Boolean-Functions-by-Ryan-ODonnell}
  {pdf-ffbe9c7b3306-p022-b004,
    pdf-ffbe9c7b3306-p023-b001}
\proofstep
  {FABL.numericalSupport}
  {context}
  {Analysis-of-Boolean-Functions-by-Ryan-ODonnell}
  {pdf-ffbe9c7b3306-p149-b001}
\proofstep
  {FABL.mem_numericalSupport}
  {context}
  {Analysis-of-Boolean-Functions-by-Ryan-ODonnell}
  {pdf-ffbe9c7b3306-p149-b001}
\proofstep
  {FABL.numericalDegree}
  {context}
  {Analysis-of-Boolean-Functions-by-Ryan-ODonnell}
  {pdf-ffbe9c7b3306-p149-b001,
    pdf-ffbe9c7b3306-p149-b004}
\proofstep
  {FABL.fourierToF₂Polynomial}
  {direct}
  {Analysis-of-Boolean-Functions-by-Ryan-ODonnell}
  {pdf-ffbe9c7b3306-p150-b001}
\proofstep
  {FABL.fourierToF₂Polynomial_one}
  {context}
  {Analysis-of-Boolean-Functions-by-Ryan-ODonnell}
  {pdf-ffbe9c7b3306-p150-b001}
\proofstep
  {FABL.PseudoBooleanFunction}
  {context}
  {Analysis-of-Boolean-Functions-by-Ryan-ODonnell}
  {pdf-ffbe9c7b3306-p149-b001}
\proofstep
  {FABL.NumericalCoefficients}
  {context}
  {Analysis-of-Boolean-Functions-by-Ryan-ODonnell}
  {pdf-ffbe9c7b3306-p149-b001}
Let $n,D\ge 0$, and let $c$ assign a real coefficient to each subset of the
$n$-coordinate set. The numerical degree of $c$ is at most $D$ if and only if every
subset with nonzero coefficient has cardinality at most $D$.
\end{theorem}

\begin{definition}[Numerical degree of a function]
\lean{FABL.functionNumericalDegree}
\statementsource{Analysis-of-Boolean-Functions-by-Ryan-ODonnell}{pdf-ffbe9c7b3306-p036-b002}
\label{FABL.functionNumericalDegree}
\leanok
\uses{FABL.PseudoBooleanFunction, FABL.numericalCoeff, FABL.numericalDegree}
For a real-valued function $\\varphi$ on $\\mathbb F_2^n$, define its numerical degree
to be the numerical degree of its numerical coefficient family.
\end{definition}

\begin{definition}[Fourier substitution coefficients]
\lean{FABL.fourierSubstitutionCoeff}
\statementsource{Analysis-of-Boolean-Functions-by-Ryan-ODonnell}{
  pdf-ffbe9c7b3306-p022-b004,
  pdf-ffbe9c7b3306-p025-b004,
  pdf-ffbe9c7b3306-p150-b001
}
\label{FABL.fourierSubstitutionCoeff}
\leanok
\uses{FABL.NumericalCoefficients, FABL.fourierCoeff}
Let $p:\\{-1,1\\}^n\\to\\mathbb R$. For each $T\\subseteq[n]$, define its substitution
coefficient by $(-2)^{|T|}\\sum_{S\\supseteq T}\\widehat p(S)$, where $\\widehat p(S)$
is the Fourier coefficient of $p$ at $S$.
\end{definition}

\begin{definition}[Fourier substitution]
\lean{FABL.fourierSubstitution}
\statementsource{Analysis-of-Boolean-Functions-by-Ryan-ODonnell}{
  pdf-ffbe9c7b3306-p022-b002,
  pdf-ffbe9c7b3306-p022-b004,
  pdf-ffbe9c7b3306-p150-b001
}
\label{FABL.fourierSubstitution}
\leanok
\uses{FABL.PseudoBooleanFunction, FABL.binaryCubeSignEquiv}
For a function $p$ on the sign cube $\\{-1,1\\}^n$, define its substitution on the
$n$-dimensional binary cube by evaluating $p$ at the sign vector corresponding to each
binary input.
\end{definition}

\begin{definition}[Fourier-to-binary function conversion]
\lean{FABL.fourierToF₂Polynomial}
\leanok
\uses{FABL.PseudoBooleanFunction, FABL.binaryCubeSignEquiv}
\statementsource{Analysis-of-Boolean-Functions-by-Ryan-ODonnell}{
  pdf-ffbe9c7b3306-p035-b010,
  pdf-ffbe9c7b3306-p036-b001,
  pdf-ffbe9c7b3306-p149-b004,
  pdf-ffbe9c7b3306-p150-b001
}
For a function $p$ on the sign cube $\\{-1,1\\}^n$, define a real-valued function on the
$n$-dimensional binary cube by assigning to each input $x$ the value $(1-p(s(x)))/2$,
where $s(x)$ is the corresponding sign vector.
\end{definition}

\begin{definition}[Fourier-to-binary coefficients]
\lean{FABL.fourierToF₂Coeff}
\leanok
\uses{FABL.NumericalCoefficients, FABL.fourierSubstitutionCoeff}
\statementsource{Analysis-of-Boolean-Functions-by-Ryan-ODonnell}{
  pdf-ffbe9c7b3306-p022-b004,
  pdf-ffbe9c7b3306-p035-b010,
  pdf-ffbe9c7b3306-p036-b001,
  pdf-ffbe9c7b3306-p150-b001
}
For a function $p$ on the sign cube $\\{-1,1\\}^n$, define a real coefficient for each
subset $T\subseteq\{1,\ldots,n\}$ as one half of the difference between the indicator
that $T$ is empty and the Fourier coefficient of $p$ at $T$.
\end{definition}

\begin{theorem}[Sign encoding of a binary coordinate]
\lean{FABL.signValue_binaryCubeSignEquiv_eq_one_sub_two}
\proofsource{Analysis-of-Boolean-Functions-by-Ryan-ODonnell}{pdf-ffbe9c7b3306-p022-b004}
\label{FABL.signValue_binaryCubeSignEquiv_eq_one_sub_two}
\leanok
\uses{FABL.F₂Cube, FABL.binaryCubeSignEquiv, FABL.binaryCubeSignEquiv_apply, FABL.numericalMonomial, FABL.signValue}
\difficulty{4}
\proofstep
  {FABL.𝔽₂}
  {context}
  {Analysis-of-Boolean-Functions-by-Ryan-ODonnell}
  {pdf-ffbe9c7b3306-p015-b008,
    pdf-ffbe9c7b3306-p148-b003}
\proofstep
  {FABL.F₂Cube}
  {context}
  {Analysis-of-Boolean-Functions-by-Ryan-ODonnell}
  {pdf-ffbe9c7b3306-p023-b001,
    pdf-ffbe9c7b3306-p148-b003}
\proofstep
  {FABL.Sign}
  {context}
  {Analysis-of-Boolean-Functions-by-Ryan-ODonnell}
  {pdf-ffbe9c7b3306-p020-b002,
    pdf-ffbe9c7b3306-p022-b004}
\proofstep
  {FABL.signValue}
  {context}
  {Analysis-of-Boolean-Functions-by-Ryan-ODonnell}
  {pdf-ffbe9c7b3306-p020-b002,
    pdf-ffbe9c7b3306-p022-b004}
\proofstep
  {FABL.signEncode}
  {direct}
  {Analysis-of-Boolean-Functions-by-Ryan-ODonnell}
  {pdf-ffbe9c7b3306-p014-b002,
    pdf-ffbe9c7b3306-p022-b004}
\proofstep
  {FABL.binarySignEquiv}
  {context}
  {Analysis-of-Boolean-Functions-by-Ryan-ODonnell}
  {pdf-ffbe9c7b3306-p020-b001,
    pdf-ffbe9c7b3306-p022-b004}
\proofstep
  {FABL.binaryCubeSignEquiv}
  {context}
  {Analysis-of-Boolean-Functions-by-Ryan-ODonnell}
  {pdf-ffbe9c7b3306-p020-b002,
    pdf-ffbe9c7b3306-p023-b001}
\proofstep
  {FABL.binaryCubeSignEquiv_apply}
  {context}
  {Analysis-of-Boolean-Functions-by-Ryan-ODonnell}
  {pdf-ffbe9c7b3306-p014-b002,
    pdf-ffbe9c7b3306-p023-b001}
\proofstep
  {FABL.numericalSupport}
  {context}
  {Analysis-of-Boolean-Functions-by-Ryan-ODonnell}
  {pdf-ffbe9c7b3306-p022-b001,
    pdf-ffbe9c7b3306-p149-b001}
\proofstep
  {FABL.mem_numericalSupport}
  {context}
  {Analysis-of-Boolean-Functions-by-Ryan-ODonnell}
  {pdf-ffbe9c7b3306-p022-b001,
    pdf-ffbe9c7b3306-p149-b001}
\proofstep
  {FABL.fourierToF₂Polynomial}
  {direct}
  {Analysis-of-Boolean-Functions-by-Ryan-ODonnell}
  {pdf-ffbe9c7b3306-p150-b001}
\proofstep
  {FABL.fourierToF₂Polynomial_one}
  {context}
  {Analysis-of-Boolean-Functions-by-Ryan-ODonnell}
  {pdf-ffbe9c7b3306-p150-b001}
\proofstep
  {FABL.PseudoBooleanFunction}
  {context}
  {Analysis-of-Boolean-Functions-by-Ryan-ODonnell}
  {pdf-ffbe9c7b3306-p148-b003,
    pdf-ffbe9c7b3306-p149-b001}
\proofstep
  {FABL.NumericalCoefficients}
  {context}
  {Analysis-of-Boolean-Functions-by-Ryan-ODonnell}
  {pdf-ffbe9c7b3306-p149-b001,
    pdf-ffbe9c7b3306-p022-b001}
\proofstep
  {FABL.numericalMonomial}
  {context}
  {Analysis-of-Boolean-Functions-by-Ryan-ODonnell}
  {pdf-ffbe9c7b3306-p148-b005,
    pdf-ffbe9c7b3306-p149-b001}
Let $x\in(\mathbb F_2)^n$ and let $i$ be a coordinate. Under the coordinatewise encoding
of bits by signs, the real value of the sign at coordinate $i$ is $1-2\mathbf
1_{\{x_i=1\}}$, where the indicator is $1$ when $x_i=1$ and $0$ otherwise.
\end{theorem}

\begin{theorem}[Binary monomial expansion]
\lean{FABL.monomial_binaryCubeSignEquiv_eq_numerical_sum}
\label{FABL.monomial_binaryCubeSignEquiv_eq_numerical_sum}
\leanok
\uses{FABL.F₂Cube, FABL.binaryCubeSignEquiv, FABL.monomial, FABL.numericalMonomial, FABL.signValue, FABL.signValue_binaryCubeSignEquiv_eq_one_sub_two}
\difficulty{4}
Let $n$ be a natural number, let $S$ be a set of coordinates in $\\{0,\\ldots,n-1\\}$,
and let $x$ assign a bit in $\\mathbb{F}_2$ to each coordinate. Encode each bit as the
sign $(-1)^b$. The sign monomial on $S$ equals $\\displaystyle\\sum_{T\\subseteq
S}(-2)^{|T|}\\prod_{i\\in T}\\mathbf{1}_{\\{x_i=1\\}}$, where the product over the empty
set is $1$.
\end{theorem}

\begin{theorem}[Evaluation of Fourier substitution coefficients]
\lean{FABL.numericalEval_fourierSubstitutionCoeff}
\label{FABL.numericalEval_fourierSubstitutionCoeff}
\leanok
\uses{FABL.binaryCubeSignEquiv, FABL.fourierCoeff, FABL.fourierSubstitution, FABL.fourierSubstitutionCoeff, FABL.fourier_expansion, FABL.monomial_binaryCubeSignEquiv_eq_numerical_sum, FABL.numericalEval, FABL.numericalMonomial}
\difficulty{5}
For any $n\ge 0$ and any function $p:\{\pm1\}^n\to\mathbb{R}$, evaluating the numerical
coefficient function $T\mapsto(-2)^{|T|}\sum_{S\supseteq T}\widehat p(S)$ in the
squarefree indicator monomials gives the function $x\mapsto p(\operatorname{sign}(x))$
on $(\mathbb{F}_2)^n$, where $\widehat p(S)$ denotes the uniform Fourier coefficient of
$p$ at $S$.
\end{theorem}

\begin{theorem}[Coefficients after Fourier substitution]
\lean{FABL.numericalCoeff_fourierSubstitution}
\label{FABL.numericalCoeff_fourierSubstitution}
\leanok
\uses{FABL.fourierSubstitution, FABL.fourierSubstitutionCoeff, FABL.numericalCoeff, FABL.numericalEval, FABL.numericalEval_fourierSubstitutionCoeff, FABL.numericalEval_injective, FABL.numericalEval_numericalCoeff}
\difficulty{3}
For any $n\ge 0$ and any function $p:\{\pm1\}^n\to\mathbb{R}$, the numerical coefficient
function of $x\mapsto p(\operatorname{sign}(x))$ is the coefficient function $c$ given,
for each coordinate subset $T$, by $\displaystyle c(T)=(-2)^{|T|}\sum_{S\supseteq
T}\widehat p(S)$, where $\widehat p(S)$ is the uniform Fourier coefficient of $p$ at
$S$.
\end{theorem}

\begin{theorem}[Fourier coefficient expansion]
\lean{FABL.numericalEval_fourierToF₂Coeff}
\leanok
\uses{FABL.binaryCubeSignEquiv, FABL.fourierSubstitutionCoeff, FABL.fourierToF₂Coeff, FABL.fourierToF₂Polynomial, FABL.numericalEval, FABL.numericalEval_fourierSubstitutionCoeff, FABL.numericalMonomial}
\difficulty{4}
For each natural number $n$ and each function $p$ from the sign cube $\\{-1,1\\}^n$ to
$\\mathbb R$, evaluate on the binary cube the polynomial whose coefficient at a subset
$T\\subseteq\\{1,\\ldots,n\\}$ is half the difference between the indicator of
$T=\\varnothing$ and $(-2)^{|T|}$ times the sum of the Fourier coefficients of $p$ over
all subsets containing $T$. This evaluation equals the function on the binary cube given
by $x\\mapsto (1-p(\\sigma(x)))/2$, where $\\sigma$ sends each binary coordinate to its
corresponding sign.
\end{theorem}

\begin{theorem}[Coefficients of the normalized Fourier transform]
\lean{FABL.numericalCoeff_fourierToF₂Polynomial}
\leanok
\uses{FABL.fourierToF₂Coeff, FABL.fourierToF₂Polynomial, FABL.numericalCoeff, FABL.numericalEval, FABL.numericalEval_fourierToF₂Coeff, FABL.numericalEval_injective, FABL.numericalEval_numericalCoeff}
\difficulty{3}
For any $n\ge 0$ and any function $p:\{\pm1\}^n\to\mathbb{R}$, the numerical coefficient
function of $x\mapsto (1-p(\operatorname{sign}(x)))/2$ has, at each coordinate subset
$T$, the value $\displaystyle
\frac{\mathbf{1}_{\{T=\varnothing\}}-(-2)^{|T|}\sum_{S\supseteq T}\widehat p(S)}{2}$,
where $\widehat p(S)$ is the uniform Fourier coefficient of $p$ at $S$.
\end{theorem}

\begin{theorem}[Fourier expansion of the transformed function]
\lean{FABL.fourierToF₂Polynomial_apply_fourierExpansion}
\leanok
\uses{FABL.F₂Cube, FABL.binaryCubeSignEquiv, FABL.fourierCoeff, FABL.fourierToF₂Polynomial, FABL.fourier_expansion, FABL.monomial_binaryCubeSignEquiv_eq_numerical_sum, FABL.numericalMonomial}
\difficulty{3}
\proofsource{Analysis-of-Boolean-Functions-by-Ryan-ODonnell}{pdf-ffbe9c7b3306-p150-b001}
\proofstep
  {FABL.𝔽₂}
  {direct}
  {Analysis-of-Boolean-Functions-by-Ryan-ODonnell}
  {pdf-ffbe9c7b3306-p015-b008}
\proofstep
  {FABL.F₂Cube}
  {context}
  {Analysis-of-Boolean-Functions-by-Ryan-ODonnell}
  {pdf-ffbe9c7b3306-p023-b002}
\proofstep
  {FABL.Sign}
  {context}
  {Analysis-of-Boolean-Functions-by-Ryan-ODonnell}
  {pdf-ffbe9c7b3306-p014-b002,
    pdf-ffbe9c7b3306-p020-b002}
\proofstep
  {FABL.signValue}
  {context}
  {Analysis-of-Boolean-Functions-by-Ryan-ODonnell}
  {pdf-ffbe9c7b3306-p014-b002}
\proofstep
  {FABL.indicatorPolynomial}
  {direct}
  {Analysis-of-Boolean-Functions-by-Ryan-ODonnell}
  {pdf-ffbe9c7b3306-p021-b002}
\proofstep
  {FABL.indicatorPolynomial_eq_ite}
  {direct}
  {Analysis-of-Boolean-Functions-by-Ryan-ODonnell}
  {pdf-ffbe9c7b3306-p021-b002}
\proofstep
  {FABL.sum_indicatorPolynomial}
  {direct}
  {Analysis-of-Boolean-Functions-by-Ryan-ODonnell}
  {pdf-ffbe9c7b3306-p021-b002}
\proofstep
  {FABL.monomial}
  {direct}
  {Analysis-of-Boolean-Functions-by-Ryan-ODonnell}
  {pdf-ffbe9c7b3306-p022-b001}
\proofstep
  {FABL.signMonomialChar}
  {context}
  {Analysis-of-Boolean-Functions-by-Ryan-ODonnell}
  {pdf-ffbe9c7b3306-p023-b001}
\proofstep
  {FABL.signMonomialChar_injective}
  {context}
  {Analysis-of-Boolean-Functions-by-Ryan-ODonnell}
  {pdf-ffbe9c7b3306-p023-b006}
\proofstep
  {FABL.expect_monomial_mul}
  {direct}
  {Analysis-of-Boolean-Functions-by-Ryan-ODonnell}
  {pdf-ffbe9c7b3306-p024-b003,
    pdf-ffbe9c7b3306-p024-b005,
    pdf-ffbe9c7b3306-p024-b007}
\proofstep
  {FABL.fourierCoeff}
  {direct}
  {Analysis-of-Boolean-Functions-by-Ryan-ODonnell}
  {pdf-ffbe9c7b3306-p025-b004}
\proofstep
  {FABL.indicatorPolynomial_fourier_sum}
  {context}
  {Analysis-of-Boolean-Functions-by-Ryan-ODonnell}
  {pdf-ffbe9c7b3306-p021-b002,
    pdf-ffbe9c7b3306-p022-b002}
\proofstep
  {FABL.fourier_expansion_from_interpolation}
  {direct}
  {Analysis-of-Boolean-Functions-by-Ryan-ODonnell}
  {pdf-ffbe9c7b3306-p022-b002}
\proofstep
  {FABL.multilinearPolynomial}
  {context}
  {Analysis-of-Boolean-Functions-by-Ryan-ODonnell}
  {pdf-ffbe9c7b3306-p020-b005,
    pdf-ffbe9c7b3306-p022-b002}
\proofstep
  {FABL.fourier_expansion_unique}
  {direct}
  {Analysis-of-Boolean-Functions-by-Ryan-ODonnell}
  {pdf-ffbe9c7b3306-p022-b002,
    pdf-ffbe9c7b3306-p023-b006}
\proofstep
  {FABL.fourier_expansion}
  {direct}
  {Analysis-of-Boolean-Functions-by-Ryan-ODonnell}
  {pdf-ffbe9c7b3306-p022-b002}
\proofstep
  {FABL.signEncode}
  {direct}
  {Analysis-of-Boolean-Functions-by-Ryan-ODonnell}
  {pdf-ffbe9c7b3306-p014-b002,
    pdf-ffbe9c7b3306-p022-b004}
\proofstep
  {FABL.binarySignEquiv}
  {context}
  {Analysis-of-Boolean-Functions-by-Ryan-ODonnell}
  {pdf-ffbe9c7b3306-p022-b004}
\proofstep
  {FABL.binaryCubeSignEquiv}
  {context}
  {Analysis-of-Boolean-Functions-by-Ryan-ODonnell}
  {pdf-ffbe9c7b3306-p022-b004}
\proofstep
  {FABL.binaryCubeSignEquiv_apply}
  {context}
  {Analysis-of-Boolean-Functions-by-Ryan-ODonnell}
  {pdf-ffbe9c7b3306-p022-b004}
\proofstep
  {FABL.numericalSupport}
  {context}
  {Analysis-of-Boolean-Functions-by-Ryan-ODonnell}
  {pdf-ffbe9c7b3306-p018-b003,
    pdf-ffbe9c7b3306-p149-b001}
\proofstep
  {FABL.mem_numericalSupport}
  {context}
  {Analysis-of-Boolean-Functions-by-Ryan-ODonnell}
  {pdf-ffbe9c7b3306-p018-b003,
    pdf-ffbe9c7b3306-p149-b001}
\proofstep
  {FABL.fourierToF₂Polynomial}
  {context}
  {Analysis-of-Boolean-Functions-by-Ryan-ODonnell}
  {pdf-ffbe9c7b3306-p150-b001}
\proofstep
  {FABL.signValue_binaryCubeSignEquiv_eq_one_sub_two}
  {context}
  {Analysis-of-Boolean-Functions-by-Ryan-ODonnell}
  {pdf-ffbe9c7b3306-p014-b002,
    pdf-ffbe9c7b3306-p022-b004,
    pdf-ffbe9c7b3306-p148-b005}
\proofstep
  {FABL.monomial_binaryCubeSignEquiv_eq_numerical_sum}
  {context}
  {Analysis-of-Boolean-Functions-by-Ryan-ODonnell}
  {pdf-ffbe9c7b3306-p150-b001,
    pdf-ffbe9c7b3306-p148-b005}
\proofstep
  {FABL.fourierToF₂Polynomial_one}
  {context}
  {Analysis-of-Boolean-Functions-by-Ryan-ODonnell}
  {pdf-ffbe9c7b3306-p150-b001}
\proofstep
  {FABL.PseudoBooleanFunction}
  {context}
  {Analysis-of-Boolean-Functions-by-Ryan-ODonnell}
  {pdf-ffbe9c7b3306-p148-b003}
\proofstep
  {FABL.NumericalCoefficients}
  {context}
  {Analysis-of-Boolean-Functions-by-Ryan-ODonnell}
  {pdf-ffbe9c7b3306-p149-b001}
\proofstep
  {FABL.numericalMonomial}
  {context}
  {Analysis-of-Boolean-Functions-by-Ryan-ODonnell}
  {pdf-ffbe9c7b3306-p148-b005,
    pdf-ffbe9c7b3306-p149-b001}
Let $p$ be any real-valued function on the $n$-dimensional sign cube, and let
$x\in\{0,1\}^n$. Write $\widehat p(S)$ for the Fourier coefficient of $p$ at $S$. Then
$\displaystyle\frac{1-p((-1)^{x_1},\ldots,(-1)^{x_n})}{2}=\frac{1-\sum_{S\subseteq\{1,\ldots,n\}}\widehat
p(S)\sum_{T\subseteq S}(-2)^{|T|}\prod_{i\in T}\mathbf{1}\{x_i=1\}}{2}$.
\end{theorem}

\begin{theorem}[A maximal-degree nonzero Fourier coefficient]
\lean{FABL.exists_top_fourierCoeff}
\label{FABL.exists_top_fourierCoeff}
\leanok
\uses{FABL.fourierCoeff, FABL.fourierDegree, FABL.fourierSupport, FABL.mem_fourierSupport}
\difficulty{4}
Let $p$ be a real-valued function on the $n$-dimensional sign cube, and suppose its
Fourier degree is nonzero. Then some set of coordinates $S$ has a nonzero Fourier
coefficient and cardinality equal to the Fourier degree of $p$.
\end{theorem}

\begin{theorem}[Fourier coefficients of supersets at maximal degree]
\lean{FABL.sum_fourierCoeff_supersets_eq_self_of_card_eq_degree}
\label{FABL.sum_fourierCoeff_supersets_eq_self_of_card_eq_degree}
\leanok
\uses{FABL.fourierCoeff, FABL.fourierDegree, FABL.fourierDegree_le_iff}
\difficulty{4}
Let $n$ be a natural number, let $p:\{-1,1\}^n\to\mathbb R$, and let
$S\subseteq\{0,\ldots,n-1\}$. If the cardinality of $S$ equals the Fourier degree of $p$
(the largest cardinality of a set with nonzero Fourier coefficient), then the sum of the
Fourier coefficients indexed by all sets $U$ containing $S$ equals the Fourier
coefficient indexed by $S$.
\end{theorem}

\begin{theorem}[Degree under Fourier substitution]
\lean{FABL.numericalDegree_fourierSubstitutionCoeff}
\label{FABL.numericalDegree_fourierSubstitutionCoeff}
\leanok
\uses{FABL.exists_top_fourierCoeff, FABL.fourierCoeff, FABL.fourierDegree, FABL.fourierDegree_le_iff, FABL.fourierSubstitutionCoeff, FABL.numericalDegree, FABL.numericalDegree_le_iff, FABL.sum_fourierCoeff_supersets_eq_self_of_card_eq_degree}
\difficulty{5}
For any $n\ge 0$ and any function $p:\{\pm1\}^n\to\mathbb{R}$, the numerical degree of
the coefficient function $T\mapsto(-2)^{|T|}\sum_{S\supseteq T}\widehat p(S)$ equals the
Fourier degree of $p$.
\end{theorem}

\begin{theorem}[Numerical degree under Fourier substitution]
\lean{FABL.functionNumericalDegree_fourierSubstitution}
\label{FABL.functionNumericalDegree_fourierSubstitution}
\leanok
\uses{FABL.fourierDegree, FABL.fourierSubstitution, FABL.functionNumericalDegree, FABL.numericalCoeff_fourierSubstitution, FABL.numericalDegree_fourierSubstitutionCoeff}
\difficulty{1}
For every nonnegative integer $n$ and every real-valued function $p$ on the sign cube
$\{ -1,1\}^n$, the numerical degree of the function obtained by substituting the
coordinatewise sign encoding of a binary input into $p$ equals the Fourier degree of
$p$.
\end{theorem}

\begin{theorem}[Degree of normalized Fourier coefficients]
\lean{FABL.numericalDegree_fourierToF₂Coeff}
\leanok
\uses{FABL.exists_top_fourierCoeff, FABL.fourierDegree, FABL.fourierSubstitutionCoeff, FABL.fourierToF₂Coeff, FABL.numericalDegree, FABL.numericalDegree_fourierSubstitutionCoeff, FABL.numericalDegree_le_iff, FABL.sum_fourierCoeff_supersets_eq_self_of_card_eq_degree}
\difficulty{5}
For any $n\ge 0$ and any function $p:\{\pm1\}^n\to\mathbb{R}$, the numerical degree of
the coefficient function
$T\mapsto\bigl(\mathbf{1}_{\{T=\varnothing\}}-(-2)^{|T|}\sum_{S\supseteq T}\widehat
p(S)\bigr)/2$ equals the Fourier degree of $p$.
\end{theorem}

\begin{theorem}[Function numerical degree fourier to f₂polynomial]
\lean{FABL.functionNumericalDegree_fourierToF₂Polynomial}
\leanok
\uses{FABL.fourierDegree, FABL.fourierToF₂Polynomial, FABL.functionNumericalDegree, FABL.numericalCoeff_fourierToF₂Polynomial, FABL.numericalDegree_fourierToF₂Coeff}
\statementsource{Analysis-of-Boolean-Functions-by-Ryan-ODonnell}{
  pdf-ffbe9c7b3306-p150-b001,
  pdf-ffbe9c7b3306-p150-b002
}
\difficulty{1}
% dataset-pipeline: stub -- prose pending build_informal_statements.py
Statement pending regeneration.
\end{theorem}

\begin{theorem}[Transform of the constant one function]
\lean{FABL.fourierToF₂Polynomial_one}
\leanok
\uses{FABL.fourierToF₂Polynomial}
\difficulty{3}
For every nonnegative integer $n$, the function on the bit cube obtained from the
constant-one function on the sign cube by the transformation $p\mapsto(1-p)/2$ is
identically zero.
\end{theorem}

\begin{definition}[Integer coefficients of a Boolean function]
\lean{FABL.booleanNumericalCoeffInt}
\statementsource{Analysis-of-Boolean-Functions-by-Ryan-ODonnell}{pdf-ffbe9c7b3306-p036-b001}
\label{FABL.booleanNumericalCoeffInt}
\leanok
\uses{FABL.F₂BooleanFunction, FABL.f₂CubeOfFinset}
Given a function $f$ from the $n$-dimensional binary cube to the two-element field and a
set $S$ of coordinates, define an integer by summing over all subsets $T\subseteq S$:
include $(-1)^{|S|-|T|}$ when $f$ is $1$ at the binary vector whose $1$-coordinates are
exactly $T$, and include $0$ otherwise.
\end{definition}

\begin{theorem}[Boolean coefficients under real embedding]
\lean{FABL.numericalCoeff_booleanRealEmbedding_eq_intCast}
\label{FABL.numericalCoeff_booleanRealEmbedding_eq_intCast}
\leanok
\uses{FABL.F₂BooleanFunction, FABL.booleanNumericalCoeffInt, FABL.booleanRealEmbedding, FABL.f₂CubeOfFinset, FABL.numericalCoeff, FABL.numericalCoeff_eq_mobius_sum}
\difficulty{4}
For any $n\ge 0$, let $f$ be a function from the binary cube $(\mathbb{F}_2)^n$ to
$\mathbb{F}_2$, and let $S\subseteq\{1,\ldots,n\}$. The numerical coefficient of the
real-valued indicator of $f$ at $S$ equals the integer alternating sum over all
$T\subseteq S$, viewed as a real number: $\displaystyle \sum_{T\subseteq
S}(-1)^{|S|-|T|}\,\mathbf{1}_{\{f(\mathbf{1}_T)=1\}}$, where $\mathbf{1}_T$ is the
binary cube point that is $1$ exactly on $T$.
\end{theorem}

\begin{theorem}[Integer coefficient modulo two]
\lean{FABL.booleanNumericalCoeffInt_cast_f₂_eq_anfCoeff}
\leanok
\uses{FABL.F₂BooleanFunction, FABL.anfCoeff, FABL.booleanNumericalCoeffInt, FABL.f₂CubeOfFinset, FABL.𝔽₂}
\difficulty{4}
Let $n$ be a natural number, let $f$ map the $n$-dimensional cube over the field with
two elements to that field, and let $S$ be a set of coordinates. The integer sum over
$T\subseteq S$ of $(-1)^{|S|-|T|}$ times the indicator that $f$ is $1$ on the vector
with ones exactly in $T$, reduced modulo two, equals the coefficient of the monomial
indexed by $S$ in the algebraic normal form of $f$.
\end{theorem}

\begin{theorem}[Transform of a Boolean sign function]
\lean{FABL.fourierToF₂Polynomial_booleanFunction}
\leanok
\uses{FABL.BooleanFunction, FABL.BooleanFunction.toReal, FABL.booleanFunctionF₂Encoding, FABL.booleanRealEmbedding, FABL.booleanRealEmbedding_booleanFunctionF₂Encoding_apply, FABL.fourierToF₂Polynomial}
\difficulty{1}
\proofsource{Analysis-of-Boolean-Functions-by-Ryan-ODonnell}{
  pdf-ffbe9c7b3306-p035-b010,
  pdf-ffbe9c7b3306-p036-b001,
  pdf-ffbe9c7b3306-p150-b001
}
\proofstep
  {FABL.𝔽₂}
  {context}
  {Analysis-of-Boolean-Functions-by-Ryan-ODonnell}
  {pdf-ffbe9c7b3306-p015-b008,
    pdf-ffbe9c7b3306-p014-b002}
\proofstep
  {FABL.F₂Cube}
  {context}
  {Analysis-of-Boolean-Functions-by-Ryan-ODonnell}
  {pdf-ffbe9c7b3306-p023-b001}
\proofstep
  {FABL.Sign}
  {context}
  {Analysis-of-Boolean-Functions-by-Ryan-ODonnell}
  {pdf-ffbe9c7b3306-p020-b001}
\proofstep
  {FABL.signValue}
  {context}
  {Analysis-of-Boolean-Functions-by-Ryan-ODonnell}
  {pdf-ffbe9c7b3306-p020-b001}
\proofstep
  {FABL.BooleanFunction}
  {direct}
  {Analysis-of-Boolean-Functions-by-Ryan-ODonnell}
  {pdf-ffbe9c7b3306-p020-b001}
\proofstep
  {FABL.BooleanFunction.toReal}
  {context}
  {Analysis-of-Boolean-Functions-by-Ryan-ODonnell}
  {pdf-ffbe9c7b3306-p020-b001}
\proofstep
  {FABL.signEncode}
  {context}
  {Analysis-of-Boolean-Functions-by-Ryan-ODonnell}
  {pdf-ffbe9c7b3306-p014-b002,
    pdf-ffbe9c7b3306-p022-b004}
\proofstep
  {FABL.binarySignEquiv}
  {context}
  {Analysis-of-Boolean-Functions-by-Ryan-ODonnell}
  {pdf-ffbe9c7b3306-p014-b002}
\proofstep
  {FABL.binaryCubeSignEquiv}
  {context}
  {Analysis-of-Boolean-Functions-by-Ryan-ODonnell}
  {pdf-ffbe9c7b3306-p022-b004}
\proofstep
  {FABL.F₂BooleanFunction}
  {context}
  {Analysis-of-Boolean-Functions-by-Ryan-ODonnell}
  {pdf-ffbe9c7b3306-p149-b001,
    pdf-ffbe9c7b3306-p149-b002}
\proofstep
  {FABL.booleanFunctionF₂Encoding}
  {context}
  {Analysis-of-Boolean-Functions-by-Ryan-ODonnell}
  {pdf-ffbe9c7b3306-p150-b001}
\proofstep
  {FABL.booleanRealEmbedding}
  {context}
  {Analysis-of-Boolean-Functions-by-Ryan-ODonnell}
  {pdf-ffbe9c7b3306-p148-b003}
\proofstep
  {FABL.booleanRealEmbedding_booleanFunctionF₂Encoding_apply}
  {context}
  {Analysis-of-Boolean-Functions-by-Ryan-ODonnell}
  {pdf-ffbe9c7b3306-p150-b001}
\proofstep
  {FABL.numericalSupport}
  {context}
  {Analysis-of-Boolean-Functions-by-Ryan-ODonnell}
  {pdf-ffbe9c7b3306-p022-b001}
\proofstep
  {FABL.mem_numericalSupport}
  {context}
  {Analysis-of-Boolean-Functions-by-Ryan-ODonnell}
  {pdf-ffbe9c7b3306-p022-b001}
\proofstep
  {FABL.fourierToF₂Polynomial}
  {context}
  {Analysis-of-Boolean-Functions-by-Ryan-ODonnell}
  {pdf-ffbe9c7b3306-p150-b001}
\proofstep
  {FABL.fourierToF₂Polynomial_one}
  {context}
  {Analysis-of-Boolean-Functions-by-Ryan-ODonnell}
  {pdf-ffbe9c7b3306-p150-b001}
\proofstep
  {FABL.PseudoBooleanFunction}
  {context}
  {Analysis-of-Boolean-Functions-by-Ryan-ODonnell}
  {pdf-ffbe9c7b3306-p148-b003,
    pdf-ffbe9c7b3306-p149-b001}
\proofstep
  {FABL.NumericalCoefficients}
  {context}
  {Analysis-of-Boolean-Functions-by-Ryan-ODonnell}
  {pdf-ffbe9c7b3306-p022-b001,
    pdf-ffbe9c7b3306-p022-b002}
Let $f$ be a sign-valued function on the $n$-dimensional sign cube, and encode signs by
bits with $+1$ represented by $0$ and $-1$ by $1$. The function on the bit cube given by
$x\mapsto(1-f((-1)^{x_1},\ldots,(-1)^{x_n}))/2$ equals the indicator that the bit
encoding of $f((-1)^{x_1},\ldots,(-1)^{x_n})$ is $1$.
\end{theorem}

\begin{theorem}[Integer formula for the transformed coefficient]
\lean{FABL.fourierToF₂Coeff_booleanFunction_eq_intCast}
\leanok
\uses{FABL.BooleanFunction, FABL.BooleanFunction.toReal, FABL.booleanFunctionF₂Encoding, FABL.booleanNumericalCoeffInt, FABL.booleanRealEmbedding, FABL.fourierToF₂Coeff, FABL.fourierToF₂Polynomial, FABL.fourierToF₂Polynomial_booleanFunction, FABL.numericalCoeff, FABL.numericalCoeff_booleanRealEmbedding_eq_intCast, FABL.numericalCoeff_fourierToF₂Polynomial}
\difficulty{3}
Let $f$ be a sign-valued function on the $n$-dimensional sign cube, and encode each
input and output sign as a bit, with $+1$ encoded by $0$ and $-1$ by $1$. For every
subset $S\subseteq\{1,\ldots,n\}$, the transformed Fourier coefficient of the
real-valued sign function associated to $f$ equals the integer, viewed as a real number,
$\displaystyle\sum_{T\subseteq S}(-1)^{|S|-|T|}\mathbf{1}\{f(\mathbf{1}_T)=1\}$, where
$\mathbf{1}_T$ is the bit vector equal to $1$ on $T$ and $0$ elsewhere, and the
condition uses the bit encoding of $f$'s output.
\end{theorem}

\begin{theorem}[Algebraic degree int cast mod two le numerical degree]
\lean{FABL.algebraicDegree_intCastModTwo_le_numericalDegree}
\statementsource{Analysis-of-Boolean-Functions-by-Ryan-ODonnell}{pdf-ffbe9c7b3306-p150-b002}
\label{FABL.algebraicDegree_intCastModTwo_le_numericalDegree}
\leanok
\uses{FABL.algebraicDegree, FABL.algebraicDegree_le_iff, FABL.numericalDegree, FABL.numericalDegree_le_iff, FABL.𝔽₂}
\difficulty{4}
% dataset-pipeline: stub -- prose pending build_informal_statements.py
Statement pending regeneration.
\end{theorem}

\begin{theorem}[Function algebraic degree le function numerical degree boolean real embedding]
\lean{FABL.functionAlgebraicDegree_le_functionNumericalDegree_booleanRealEmbedding}
\statementsource{Analysis-of-Boolean-Functions-by-Ryan-ODonnell}{
  pdf-ffbe9c7b3306-p150-b001,
  pdf-ffbe9c7b3306-p150-b002
}
\label{FABL.functionAlgebraicDegree_le_functionNumericalDegree_booleanRealEmbedding}
\leanok
\uses{FABL.F₂BooleanFunction, FABL.algebraicDegree, FABL.algebraicDegree_intCastModTwo_le_numericalDegree, FABL.anfCoeff, FABL.booleanNumericalCoeffInt, FABL.booleanNumericalCoeffInt_cast_f₂_eq_anfCoeff, FABL.booleanRealEmbedding, FABL.functionAlgebraicDegree, FABL.functionNumericalDegree, FABL.numericalCoeff, FABL.numericalCoeff_booleanRealEmbedding_eq_intCast, FABL.numericalDegree, FABL.𝔽₂}
\difficulty{3}
% dataset-pipeline: stub -- prose pending build_informal_statements.py
Statement pending regeneration.
\end{theorem}

\begin{theorem}[Function algebraic degree boolean function f₂encoding le fourier degree via polynomial]
\lean{FABL.functionAlgebraicDegree_booleanFunctionF₂Encoding_le_fourierDegree_viaPolynomial}
\leanok
\uses{FABL.BooleanFunction, FABL.BooleanFunction.toReal, FABL.booleanFunctionF₂Encoding, FABL.booleanRealEmbedding, FABL.fourierDegree, FABL.fourierToF₂Polynomial, FABL.fourierToF₂Polynomial_booleanFunction, FABL.functionAlgebraicDegree, FABL.functionAlgebraicDegree_le_functionNumericalDegree_booleanRealEmbedding, FABL.functionNumericalDegree, FABL.functionNumericalDegree_fourierToF₂Polynomial}
\statementsource{Analysis-of-Boolean-Functions-by-Ryan-ODonnell}{
  pdf-ffbe9c7b3306-p150-b001,
  pdf-ffbe9c7b3306-p150-b002
}
\difficulty{3}
% dataset-pipeline: stub -- prose pending build_informal_statements.py
Statement pending regeneration.
\end{theorem}
